\documentclass[12pt,a4paper]{article}

\usepackage{fontspec}
\usepackage{geometry}
\usepackage{url}
\usepackage{hyperref}
\usepackage{float}
\usepackage{tikz}
\usepackage{listings}
\usepackage{xcolor}
\usetikzlibrary{shapes,arrows,positioning,fit,backgrounds}

\lstset{
    basicstyle=\ttfamily\small,
    breaklines=true,
    frame=single,
    numbers=left,
    numberstyle=\tiny,
    tabsize=2
}

\geometry{margin=1in}

\title{Distributed Compute Automata: A Standard for Logic and Data Privacy for Computations *draft}
\author{Osoro Bironga, fanosoro@gmail.com}
\date{\today}

\makeatletter
\renewcommand{\@maketitle}{%
  \newpage
  \null
  \vskip 2em%
  \begin{center}%
  \let \footnote \thanks
    {\LARGE \bfseries \@title \par}%
    \vskip 1.5em%
    {\large
      \lineskip .5em%
      \begin{tabular}[t]{c}%
        \@author
      \end{tabular}\par}%
    \vskip 1em%
    {\large \@date}%
  \end{center}%
  \par
  \vskip 1.5em}
\makeatother

\begin{document}

\maketitle

\begin{abstract}
\centering
\begin{minipage}{0.8\textwidth}
\centering
\textbf{Distributed Compute Automata} (DISCA) is a computation approach implemented through a Virtual Machine (VM) that uses Fully Homomorphic Encryption (FHE) and operates as a distributed computer.
DISCA is not a blockchain network; it is a distributed computing system with no dependence on any prior protocol.

Privacy-preserving computation aims to protect the privacy of data and the computation itself. This means we don't know what the data is or what the computation is doing. 

This is a very useful abstraction that can be applied to a scenario like running a proprietary algorithm in the open web without exposing the algorithm and without users exposing their data to you.

In this work, we outline DISCA as a standard for logic and data privacy for computations.
DISCA operates as a distributed computer with its own protocol, designed through a VM that uses FHE with the aspirational goal of low barrier of entry and ubiquitous deployment on diverse devices including smartphones, though current computational requirements limit practical deployment on resource-constrained devices.
Its design is proposed and we analyze security properties from a simple prototype.
Lastly, we create a formal verification proof for the distributed computer.
\end{minipage}
\end{abstract}

\noindent\textit{Note: Throughout this document, we use DISCA as the abbreviation for Distributed Compute Automata.}

\section{Introduction}

DISCA is implemented through a Virtual Machine (VM) that uses Fully Homomorphic Encryption (FHE) and operates as a distributed computer.
It is not a blockchain network, nor does it depend on any prior protocol.
DISCA functions as a distributed computing system where nodes participate in computations through a VM that uses FHE, enabling computations on encrypted data without decryption.

This work represents due diligence for larger work, consisting of informal experiments by the author.
These experiments serve to establish that there is groundwork for the approach, demonstrating feasibility and exploring the design space rather than claiming to solve all challenges in the field.

\subsection{Motivation}

There is active work and research in the field~\cite{zama_fhe,phala_network}.
DISCA focuses on Fully Homomorphic Encryption (FHE) as the primary approach.
By providing a unified interface for FHE-based computations, DISCA aims to serve as a platform for:
\begin{itemize}
    \item \textbf{Research Exploration:} Testing new FHE techniques and approaches
    \item \textbf{Practical Experimentation:} Trying out FHE-based approaches in various scenarios
    \item \textbf{Algorithm Development:} Enabling developers to create and deploy proprietary algorithms using FHE without exposing the algorithm logic
\end{itemize}

This work represents a first of many steps, rather than a definitive solution.

\subsection{Preliminaries}

Methods through software are achieved through cryptography.
More specifically we have Homomorphic Encryption (HE) and Fully Homomorphic Encryption (FHE).
There are also cryptographic hybrid solutions that combine multiple techniques, such as Homomorphic Encryption (HE) and Zero-Knowledge Proofs (ZKPs).

However, it is important to note that the technology will be truly achieved when both cost and performance graphs trend downward, enabling practical deployment at scale.
Currently, this goal remains aspirational: the necessary hardware infrastructure is not yet widely available, and software-based approaches remain computationally expensive by orders of magnitude compared to non-privacy-preserving alternatives.
DISCA serves as an experimental platform to explore these technologies as they mature, rather than claiming to solve the performance and cost challenges today.

\section{DISCA: A Distributed Computer}

DISCA is a distributed computer that implements computation through a Virtual Machine (VM) that uses Fully Homomorphic Encryption (FHE).
It is not a blockchain network, nor does it depend on any prior protocol or infrastructure.
DISCA operates as an independent distributed computing system where nodes participate in computations.

The system consists of nodes that communicate using DISCA's own protocol, operating as a distributed computer with no dependence on blockchain, smart contracts, or any other prior protocol infrastructure.
DISCA is designed with the aspirational goal of low barrier of entry, enabling nodes to eventually run on diverse devices including smartphones, making ubiquitous deployment a long-term objective of the system, though current computational requirements limit practical deployment on resource-constrained devices.

DISCA implements its own networking layer for peer-to-peer communication, network discovery, and message routing.
This independent networking implementation allows DISCA to operate as a self-contained distributed computer focused on computation through a VM that uses FHE, without dependencies on external protocols or infrastructure.

\subsection{Circuit Design}

DISCA transforms programs into FHE circuits through a systematic approach that leverages WebAssembly (WASM) as an intermediate representation.
This approach enables the construction of circuits from standard programming languages while maintaining the security and isolation properties required for distributed execution.

\subsubsection{Compilation to WebAssembly}

Programs written in high-level languages (Rust, C/C++, Go, etc.) are first compiled to WebAssembly format.
This compilation step is crucial for several reasons:

\begin{itemize}
    \item \textbf{WASM Stack Machine Standard:} We use the WASM stack machine standard to create our FHE circuits.
    The stack-based architecture provides a clean, predictable execution model that maps well to homomorphic encryption operations.
    Each stack operation can be directly translated into FHE circuit gates.

    \item \textbf{Simple Register Model:} WASM's stack-based execution model is simpler than register-based architectures, making it easier to reason about and transform into FHE circuits.
    The explicit stack operations eliminate the need to track register state across operations.

    \item \textbf{Language Support:} WASM is supported by many programming languages, allowing developers to write programs in their preferred language and compile to a common format.
    This enables a diverse ecosystem of tools and languages while maintaining a unified circuit generation pipeline.

    \item \textbf{Sandboxing:} WASM provides built-in sandboxing capabilities, ensuring that programs execute in an isolated environment with controlled memory access and execution boundaries.
    This isolation is essential for maintaining privacy guarantees in the distributed computer.
\end{itemize}

\subsubsection{Conversion to WebAssembly Text Format}

After compilation, the binary WASM file is converted to WebAssembly Text Format (WAT) for analysis and circuit generation:

\begin{verbatim}
wasm-tools print sample_program.wasm -o sample_program.wat
\end{verbatim}

The WAT format serves several critical purposes in the circuit design process:

\begin{itemize}
    \item \textbf{Step-through Register Operations:} WAT format allows us to step through stack operations instead of dealing with binary data, making it easier to understand and analyze the program's execution flow.
    Each instruction can be examined individually to determine its corresponding FHE circuit representation.

    \item \textbf{Human-Readable Analysis:} The text format enables developers and researchers to inspect the program structure, function signatures, and instruction sequences.
    This transparency facilitates verification and debugging of the circuit generation process.

    \item \textbf{Circuit Analysis:} The WAT representation makes it straightforward to analyze how operations map to FHE circuit construction.
    Each WASM instruction (e.g., \texttt{i32.add}, \texttt{i32.mul}, \texttt{i32.gt\_s}) can be mapped to corresponding FHE operations, allowing systematic circuit generation.

    \item \textbf{Debugging and Verification:} WAT format facilitates debugging and verification of the program's transformation from source code to WASM instructions, and ultimately to FHE circuits.
\end{itemize}

\subsubsection{Sample Program}

To demonstrate the circuit design approach, we provide an oversimplified example program that implements basic arithmetic operations.
The sample program includes three simple functions: \texttt{add}, \texttt{multiply}, and \texttt{subtract}.

This oversimplified example demonstrates how basic arithmetic operations are transformed into stack-based instructions that can be mapped to FHE circuit gates.
Each function is compiled to WASM and converted to WAT, providing a clear mapping from high-level operations to stack-based instructions.
Note that this example is intentionally simplified for the purpose of this paper; typical programs are much more complex, involving control flow, state management, loops, conditionals, and multi-step computations with intermediate values.

The Rust source code for the sample program is shown below:

\begin{lstlisting}[language=Rust, caption={Sample Program Source Code (lib.rs)}]
#[unsafe(no_mangle)]
pub extern "C" fn add(a: i32, b: i32) -> i32 {
    a + b
}

#[unsafe(no_mangle)]
pub extern "C" fn subtract(a: i32, b: i32) -> i32 {
    a - b
}

#[unsafe(no_mangle)]
pub extern "C" fn multiply(a: i32, b: i32) -> i32 {
    a * b
}
\end{lstlisting}

After compilation to WASM and conversion to WAT, the resulting WebAssembly Text Format representation demonstrates how high-level operations are transformed into stack-based instructions.
Note that compiler-generated WAT includes metadata such as memory management, stack pointers, and producer information.
If we remove this metadata, WebAssembly Text Format can be written by hand quite simply.
For example, a minimal hello world program can be written as:

\begin{lstlisting}[language=Lisp, caption={Minimal WebAssembly Text Format Example}]
(module
  (func (result i32)
    (i32.const 42)
  )
  (export "helloWorld" (func 0))
)
\end{lstlisting}

For our arithmetic example, removing the compiler-generated metadata yields a simplified representation:

\begin{lstlisting}[language=Lisp, caption={Simplified WebAssembly Text Format (simple\_arithmetic.wat)}]
(module
  (func $add (param i32 i32) (result i32)
    local.get 1
    local.get 0
    i32.add
  )
  (func $multiply (param i32 i32) (result i32)
    local.get 1
    local.get 0
    i32.mul
  )
  (func $subtract (param i32 i32) (result i32)
    local.get 0
    local.get 1
    i32.sub
  )
  (export "add" (func $add))
  (export "multiply" (func $multiply))
  (export "subtract" (func $subtract))
)
\end{lstlisting}

After generation of the final WAT file, the WebAssembly Standard, the text format in particular, is effectively our Intermediate Representation (IR).
We will not be using WebAssembly directly, but rather its instruction set as our IR.
This IR can now be decomposed into FHE circuits, where each WebAssembly instruction maps to corresponding homomorphic operations.

The WAT representation clearly shows how each function is decomposed into stack operations.
For example, the \texttt{add} function uses \texttt{local.get} instructions to push parameters onto the stack, followed by \texttt{i32.add} to perform the addition.
Similarly, \texttt{subtract} and \texttt{multiply} follow the same pattern with their respective operations (\texttt{i32.sub}, \texttt{i32.mul}).
These stack-based instructions provide a clear foundation for mapping to FHE circuit gates, where each operation can be represented as a homomorphic computation over encrypted values.

\section{References}

\begin{thebibliography}{9}

\bibitem{homomorphic_encryption}
Gentry, C. (2009). Fully homomorphic encryption using ideal lattices. In \textit{Proceedings of the 41st annual ACM symposium on Theory of computing} (pp. 169-178). ACM.

\bibitem{tfhe_scheme}
Chillotti, I., Gama, N., Georgieva, M., \& Izabach\`ene, M. (2020). TFHE: Fast Fully Homomorphic Encryption over the Torus. \textit{Journal of Cryptology}, 33(1), 34-91.

\bibitem{tfhe_rs}
Zama. \textit{tfhe-rs: A Pure Rust Implementation of the TFHE Scheme}. Retrieved from \url{https://github.com/zama-ai/tfhe-rs}. All evaluation measurements in this paper were taken against version 1.5.0 with default parameters.

\bibitem{zama_fhe}
Zama. (2024). Concrete: Fully Homomorphic Encryption (FHE) library in Rust. \textit{Zama}. Retrieved from \url{https://www.zama.ai/}

\bibitem{deterministic_tfhe}
Smart, N. P., \& Walter, M. (2025). Reactive Correctness, sINDCPA-D-Security and Deterministic Evaluation for TFHE. \textit{IACR Cryptology ePrint Archive}, Report 2025/2005. Retrieved from \url{https://eprint.iacr.org/2025/2005}. Shows that sINDCPA-D security ordinarily calls for a randomised evaluation procedure, and that the randomisation can be removed for TFHE in the random oracle model --- the security counterpart to the practical reproducibility requirement discussed in Section 3.2.

\bibitem{byzantine_generals}
Lamport, L., Shostak, R., \& Pease, M. (1982). The Byzantine Generals Problem. \textit{ACM Transactions on Programming Languages and Systems}, 4(3), 382-401.

\bibitem{shamir_secret_sharing}
Shamir, A. (1979). How to share a secret. \textit{Communications of the ACM}, 22(11), 612-613.

\bibitem{floating_point}
Goldberg, D. (1991). What Every Computer Scientist Should Know About Floating-Point Arithmetic. \textit{ACM Computing Surveys}, 23(1), 5-48.

\bibitem{cooley_tukey}
Cooley, J. W., \& Tukey, J. W. (1965). An algorithm for the machine calculation of complex Fourier series. \textit{Mathematics of Computation}, 19(90), 297-301.

\bibitem{wasm_spec}
World Wide Web Consortium. \textit{WebAssembly Core Specification}. Retrieved from \url{https://webassembly.github.io/spec/core/}

\bibitem{phala_network}
Phala Network. (2024). Phala Network: Privacy-preserving computation for Web3. \textit{Phala Network}. Retrieved from \url{https://phala.network/}. Phala provides cryptographic attestations that allow verification of the integrity of AI workloads in real-time, ensuring they run in genuine secure environments.

\end{thebibliography}


\end{document}